\documentclass[tikz,border=10pt,convert={density=600,outext=.png}]{standalone}
\usepackage{tikz}
\usepackage{xcolor}
\usetikzlibrary{
    shapes.geometric,
    arrows.meta,
    positioning,
    fit,
    backgrounds,
    calc,
    shadows.blur
}

\definecolor{mdgreen}{HTML}{4CAF50}
\definecolor{typstblue}{HTML}{2196F3}
\definecolor{pandocorange}{HTML}{FF9800}
\definecolor{docxpurple}{HTML}{9C27B0}
\definecolor{mergepink}{HTML}{E91E63}

\begin{document}
\begin{tikzpicture}[
    font=\sffamily\small,
    box/.style={
        rectangle,
        rounded corners=4pt,
        draw=#1!60!black,
        fill=#1!8,
        text=black,
        align=center,
        minimum width=3.2cm,
        inner sep=8pt,
        font=\small
    },
    file/.style={
        rectangle,
        rounded corners=2pt,
        draw=gray!40,
        fill=gray!3,
        align=center,
        minimum width=3.8cm,
        minimum height=0.7cm,
        font=\footnotesize\ttfamily
    },
    arrow/.style={
        -{Stealth[length=2.5pt, width=2pt]},
        thick,
        draw=gray!50,
        shorten >=1pt,
        shorten <=1pt
    },
    label/.style={
        font=\sffamily\footnotesize\bfseries,
        text=gray!60!black
    }
]

% ============================================
% Row 1: Source
% ============================================
\node[file, fill=green!5, draw=mdgreen!40] (md) {
    \textbf{markdown/*.md}\\
    chapter-*.md\quad abstract-*.md\\
    acknowledgment.md
};

\node[label, above=0.3cm of md] (srclabel) {写作源（纯文本）};

% ============================================
% Row 2: md2typst
% ============================================
\node[box={mdgreen}, below=1.2cm of md] (md2typst) {
    \textbf{md2typst.py}\\
    Markdown $\to$ Typst\\
    章节编号 + 图表标签
};

\draw[arrow] (md) -- (md2typst);

% ============================================
% Row 3: main.typ
% ============================================
\node[file, fill=blue!5, draw=typstblue!40, below=1.2cm of md2typst] (typst) {
    \textbf{src/main.typ}\\
    header.typ + \#bibliography()
};

\draw[arrow] (md2typst) -- (typst);

% ============================================
% Row 4: Pandoc
% ============================================
\node[box={typstblue}, below=1.2cm of typst, minimum width=6cm] (pandoc) {
    \textbf{pandoc}\\
    Lua filters (fix-ref + fix-table-caption + fix-pagebrk)\\
    --citeproc --csl=ieee.csl --reference-doc=clean\_ref.docx
};

\draw[arrow] (typst) -- (pandoc);

% ============================================
% Row 5: Body + TOC
% ============================================
\node[file, fill=orange!5, draw=pandocorange!40, below=1.2cm of pandoc, xshift=-2.8cm] (body) {
    \textbf{main\_content.docx}\\
    正文排版
};

\node[file, fill=pandocorange!5, draw=pandocorange!40, below=1.2cm of pandoc, xshift=2.8cm] (extended) {
    \textbf{extended\_content.docx}\\
    正文 + 目录
};

\draw[arrow] (pandoc.south) -- ++(0,-0.3) -| (body.north);
\draw[arrow] (pandoc.south) -- ++(0,-0.3) -| (extended.north);

\node[box={pandocorange}, right=0.3cm of extended, minimum width=2.2cm, minimum height=1cm] (toc) {
    \textbf{gen-toc.py}\\
    目录注入
};

\draw[arrow] (body.east) -- (toc.west);
\draw[arrow] (toc.north) -- ++(0,0.3) -| (extended.north east);

% ============================================
% Row 6: Cover (two templates)
% ============================================
\node[file, fill=purple!5, draw=docxpurple!40, below=1.8cm of body, xshift=-2.8cm] (prelimcn) {
    \textbf{prelim-cn.docx}\\
    中文封面 + 原创声明 + 摘要
};

\node[file, fill=docxpurple!5, draw=docxpurple!40, below=1.8cm of body, xshift=2.8cm] (prelimen) {
    \textbf{prelim-en.docx}\\
    English cover + declaration + abstract
};

% prelim templates
\node[file, fill=gray!3, draw=gray!30, above=0.8cm of prelimcn, minimum width=4.2cm] (prelimcn_tpl) {
    \textbf{template/prelim-cn.docx}\\
    \{\{TITLE\_CN\}\}\quad \{\{SID\}\}\quad ... \quad\{\{ABSTRACTCN\}\}
};

\node[file, fill=gray!3, draw=gray!30, above=0.8cm of prelimen, minimum width=4.2cm] (prelimen_tpl) {
    \textbf{template/prelim-en.docx}\\
    \{\{TITLE\_EN\}\}\quad \{\{SID\}\}\quad ... \quad\{\{ABSTRACTEN\}\}
};

\node[box={docxpurple}, above=0.3cm of prelimcn_tpl, xshift=1.5cm, minimum width=2cm, minimum height=0.7cm] (render) {
    \textbf{substitute-txt.py}\\
    config.json
};

\draw[arrow] (prelimcn_tpl.east) -- ++(1.2,0) |- (render.south west);
\draw[arrow] (prelimen_tpl.west) -- ++(-1.2,0) |- (render.south east);
\draw[arrow] (render) |- (prelimcn_tpl.north);
\draw[arrow] (render) |- (prelimen_tpl.north);
\draw[arrow] (prelimcn_tpl) -- (prelimcn);
\draw[arrow] (prelimen_tpl) -- (prelimen);

% ============================================
% Row 7: Merge
% ============================================
\node[box={mergepink}, below=2.2cm of prelimcn, xshift=2.8cm, minimum width=8cm, minimum height=1.4cm] (merge) {
    \textbf{merge-docs.py + setup-thesis.py}\\
    封面 $\oplus$ 正文 $\to$ 分节符合并 $\to$ 页码设置 $\to$ 页眉清除
};

\draw[arrow] (extended.south) -- ++(0,-0.5) -| ([xshift=1.5cm]merge.north);
\draw[arrow] (prelimcn.south) -- ++(0,-0.3) -| ([xshift=-1.5cm]merge.north);
\draw[arrow] (prelimen.south) -- ++(0,-0.3) -| (merge.north);

% ============================================
% Row 8: Final outputs
% ============================================
\node[file, fill=red!5, draw=mergepink!50, below=1.2cm of merge, xshift=-2cm] (finalcn) {
    \textbf{result/thesis-cn.docx}\\
    中文学位论文
};

\node[file, fill=red!5, draw=mergepink!50, below=1.2cm of merge, xshift=2cm] (finalen) {
    \textbf{result/thesis-en.docx}\\
    English thesis
};

\draw[arrow] (merge.south) -- ++(0,-0.3) -| (finalcn.north);
\draw[arrow] (merge.south) -- ++(0,-0.3) -| (finalen.north);

% ============================================
% Makefile orchestration label
% ============================================
\node[
    rectangle,
    rounded corners=3pt,
    draw=gray!30,
    fill=gray!2,
    inner sep=4pt,
    font=\sffamily\footnotesize\bfseries,
    text=gray!50,
    below=0.5cm of finalcn,
    xshift=2cm,
    minimum width=10cm
] (make) {
    \$ make all \quad$\to$\quad 一条命令，完整构建
};

\end{tikzpicture}
\end{document}
